We need two things from a simulator: that it produce paths of the process we specified,
and that we be able to tell whether it did.
This section gives one general algorithm, one exact algorithm for the exponential kernel,
and the test.

\subsection{Thinning}

\cite{LS79sim} first described an algorithm for the simulation of non-homogeneous Poisson
processes based on thinning.
\cite{Oga81lew} generalised it, replacing the requirement of bounded intensities by local
boundedness.
We recall the algorithm, since it applies to any kernel and is the control against which the
exact scheme below is checked.

The idea is Proposition \ref{prop.conditionalSurvival} run backwards.
If $\totalIntensity$ were constant, the next inter-arrival would be exponential of that rate.
It is not constant, but if it is non-increasing between events then its value at the last event
dominates it,
so one may draw an exponential at the dominating rate and keep the candidate with probability
equal to the ratio of the true rate to the dominating one.
That $\totalIntensity$ is non-increasing between events is exactly the non-negativity of the
kernel, $\excitation \geq 0$: excitation only ever decays until the next event adds to it.
The bound is therefore free, and no separate majorant has to be constructed.

We take $(\arrivalTimes[]_n,\event[n])$ as the sequence of arrival times and event labels
of the ground process, and recall that computing $\totalIntensity(t)$ requires
$\lbrace(\arrivalTimes[]_n,\event[n]): \, \arrivalTimes[]_n <t\rbrace$.

\begin{algorithm}[h!]
 \caption{{Ogata's thinning algorithm \citep{Oga81lew}}}
 \label{algo.ogata}
 \begin{algorithmic}[5]
  \REQUIRE $(\arrivalTimes[]_i, \event[i])_{i=1,\dots,n-1}$
  \STATE set $t:=\arrivalTimes[]_{n-1}$
  \STATE set $\xi:=0$
  \WHILE{$\xi=0$}
  \STATE draw $U \sim \text{Exp}(\totalIntensity(t))$
  \STATE set $\xi :=1$ with probability $\totalIntensity(t+U)/\totalIntensity(t)$
  \STATE update $t\leftarrow t+U$
  \ENDWHILE
  \STATE set $\arrivalTimes[]_n:=t$
  \STATE draw $\event[n]$ in $\lbrace 1,\dots,\numEventTypes\rbrace$ with probabilities proportional to $\lbrace \intensity[1](\arrivalTimes[]_n), \dots, \intensity[\numEventTypes](\arrivalTimes[]_n)\rbrace$
  \RETURN $(\arrivalTimes[]_n,\event[n])$
 \end{algorithmic}
\end{algorithm}

The algorithm is exact but wasteful: the rejected candidates are work done and discarded,
and the rejection rate grows with the branching ratio.
For the kernel of Section \ref{sec.exponentialKernel} it can be avoided entirely.

\subsection{An exact scheme for the exponential kernel}

Write $\totalIntensity = \mathbf{1}^{\transpose}\intensity[ ]
= \bar{\baseIntensity} + (\mathbf{1}^{\transpose}\excitation) \decayedCounts$.
Between events every component of $\decayedCounts$ decays by the same factor,
because the decay $\decay$ is common to all pairs;
hence so does the excited part of $\totalIntensity$, and
\begin{equation}\label{eq.totalIntensityBetweenEvents}
 \totalIntensity(t) = \bar{\baseIntensity}
 + D \, e^{-\decay(t - \arrivalTimes[]_n)},
 \qquad
 D := \totalIntensity(\arrivalTimes[]_n +) - \bar{\baseIntensity},
 \qquad t \in (\arrivalTimes[]_n, \arrivalTimes[]_{n+1}].
\end{equation}
It is the common $\decay$ that does this, and nothing else:
constant column sums are not needed.
Non-negativity of $\excitation$ enters separately, and gives $D \geq 0$.

Substituting \eqref{eq.totalIntensityBetweenEvents} into
Proposition \ref{prop.conditionalSurvival} and integrating,
\begin{equation}\label{eq.dzSurvival}
 \Prob \left( \arrivalTimes[]_{n+1} - \arrivalTimes[]_{n} > s \, \vert \,
 \filtrationF_{\arrivalTimes[]_{n}} \right)
 = e^{-\bar{\baseIntensity}s} \cdot
 \exp\left( -\frac{D}{\decay}\big(1 - e^{-\decay s}\big) \right).
\end{equation}
The right-hand side is a product of two survival functions,
so the inter-arrival time is the minimum of two independent random variables,
one for the immigration and one for the excitation.
This is the observation of \cite{DZ13exa}.

\begin{prop}[\cite{DZ13exa}]\label{prop.dassiosZhao}
 Let $U_1, U_2$ be independent and uniform on $(0,1)$, and set
 \begin{equation}\label{eq.dzDraws}
  \tau_1 := -\frac{\log U_1}{\bar{\baseIntensity}},
  \qquad
  \tau_2 := -\frac{1}{\decay}
  \log\left( 1 + \frac{\decay \log U_2}{D} \right),
 \end{equation}
 with $\tau_2 := +\infty$ when $1 + \decay \log U_2 / D \leq 0$.
 Then $\arrivalTimes[]_{n+1} - \arrivalTimes[]_{n}$ has the law \eqref{eq.dzSurvival}
 and is distributed as $\tau_1 \wedge \tau_2$.
 The type is then drawn at the arrival time, before the jump, with
 \begin{equation}\label{eq.dzType}
  \Prob \big( \event[n+1] = e \big)
  = \frac{\intensity(\arrivalTimes[]_{n+1}-)}{\totalIntensity(\arrivalTimes[]_{n+1}-)} .
 \end{equation}
\end{prop}

Nothing here is approximate and nothing is rejected.
Two points deserve comment.

The case $\tau_2 = +\infty$ is not a numerical edge case to be patched.
The excited part of \eqref{eq.dzSurvival} carries the finite total mass $D/\decay$,
which is the expected number of offspring the current state still owes;
with probability $e^{-D/\decay}$ it expires without ever firing,
and the next event is an immigrant.
An implementation that clips instead of returning $+\infty$ is biasing the immigrant share.

And the type draw \eqref{eq.dzType} is part of the algorithm, not an afterthought.
The waiting time \eqref{eq.dzDraws} sees the excitation matrix only through its column sums,
so it is \eqref{eq.dzType} that makes the scheme multivariate,
and it is there that the full matrix $\excitation$ re-enters.
For the same reason $\totalIntensity$ is not itself a one-dimensional Hawkes process:
it is exponential \emph{between} events, but the size of its jump depends on the type that
fires, so the full state $\decayedCounts$ must still be carried.

\subsection{Goodness of fit}

Theorem \ref{thm.meyer1971} turns the compensator into a test.
If the path is a path of the model, then rescaling each component's clock by its own
compensator turns it into a unit-rate Poisson process,
and the rescaled inter-arrival times \eqref{eq.time-changed_intertimes} are independent
$\text{Exp}(1)$ variables.
Both ingredients are available in closed form:
the compensator by \eqref{eq.compensatorClosedForm},
and the arrival times by construction.
The test is due to \cite{Oga88sta}, who introduced it for earthquake catalogues.

The three hypotheses are those of Theorem \ref{thm.meyer1971},
and it is worth restating them here rather than assuming that
Section \ref{sec.countingProcesses} is remembered:
no two components jump simultaneously;
the compensator is continuous;
and it diverges.
The second and third hold by construction for the model of this chapter.
The first holds for a simulated path, where the arrival times are almost surely distinct,
and is precisely what fails for a recorded one.

Finally, an honest statement of what such a test does and does not check.
Comparing each component's rescaled inter-arrival times against $\text{Exp}(1)$ ---
by a Kolmogorov--Smirnov statistic, say --- tests the \emph{marginal} law and nothing more.
Theorem \ref{thm.meyer1971} asserts two independences beyond that law:
serial independence within a component, and joint independence across components.
Neither is tested by a comparison of margins,
and the second is exactly what a violation of the no-common-jumps hypothesis would break.
A suite that reports six passing KS tests has established six margins.
